\uuid{7cH5}
\titre{Extrema d’une fonction sur un domaine fermé }

\niveau{}
\module{ Fonctions de plusieurs variables }
\chapitre{Optimisation sous contrainte}
\sousChapitre{}
\theme{}
\auteur{Maxime Nguyen}
\datecreate{ 2026-04-07}
\organisation{AMSCC}
\difficulte{3}

\contenu{

\texte{ 
On considère la fonction $f(x,y)=x^2+2y
$
définie sur l'ensemble
$$
K=\{(x,y)\in\mathbb{R}^2 \mid x^2+y^2\le 4,\ y\ge x\}.$$
On note :
\[
B_1=\{(x,x)\in\mathbb{R}^2 \mid -\sqrt{2}\le x\le \sqrt{2}\},\,B_2=\{(2\cos\theta,2\sin\theta)\in\mathbb{R}^2 \mid \theta\in[\tfrac{\pi}{4},\tfrac{5\pi}{4}]\}.
\]
}

\begin{enumerate}

\item   \question{Représenter graphiquement le domaine $K$.}
\indication{}
\reponse{
L’ensemble $K$ est l’intersection :
\[
K=\{(x,y)\in\mathbb{R}^2\mid x^2+y^2\le 4\}\cap \{(x,y)\in\mathbb{R}^2\mid y\ge x\}.
\]
Il s’agit donc du disque fermé de centre $(0,0)$ et de rayon $2$, coupé par le demi-plan situé au-dessus de la droite $y=x$.

Le bord de $K$ est formé de deux parties :
\begin{itemize}
\item le segment
\[
B_1=\{(x,x)\mid -\sqrt{2}\le x\le \sqrt{2}\},
\]
qui correspond à la portion de la droite $y=x$ située dans le disque ;
\item l’arc de cercle
\[
B_2=\{(2\cos\theta,2\sin\theta)\mid \theta\in[\tfrac{\pi}{4},\tfrac{5\pi}{4}]\},
\]
qui relie les points $(\sqrt{2},\sqrt{2})$ et $(-\sqrt{2},-\sqrt{2})$ en passant par la partie du cercle située dans le demi-plan $y\ge x$.
\end{itemize}

Graphiquement, $K$ est donc une moitié de disque fermé, bordée par le segment de droite $B_1$ et l’arc de cercle $B_2$.
}

%\item   \question{Justifier que $K$ est fermé et borné.}
%\indication{}
%\reponse{}

\item   \question{Montrer que $f$ admet un minimum global et un maximum global sur $K$.}
\indication{}
\reponse{
La fonction
\[
f(x,y)=x^2+2y
\]
est polynomiale, donc continue sur $\mathbb{R}^2$.

D’après la représentation précédente, $K$ est fermé et borné, donc compact.
Par le théorème de Weierstrass, toute fonction continue sur un compact admet un minimum global et un maximum global.

Ainsi, $f$ admet un minimum global et un maximum global sur $K$.
}

\item   \question{Déterminer les points critiques éventuels de $f$ dans l’intérieur de $K$.}
\indication{}
\reponse{
On calcule le gradient de $f$ :
\[
\nabla f(x,y)=\left(\frac{\partial f}{\partial x}(x,y),\frac{\partial f}{\partial y}(x,y)\right)=(2x,2).
\]
Un point critique intérieur devrait vérifier
\[
2x=0
\qquad\text{et}\qquad
2=0.
\]
La seconde équation est impossible. Il n’existe donc aucun point critique dans l’intérieur de $K$.
}

\item   \question{Étudier la restriction de $f$ à $B_1$. On posera
$g(x)=f(x,x).
$
Déterminer les variations de $g$ sur $[-\sqrt{2},\sqrt{2}]$.}
\indication{}
\reponse{
Sur $B_1$, on a $y=x$, donc
\[
g(x)=f(x,x)=x^2+2x,
\qquad x\in[-\sqrt{2},\sqrt{2}].
\]
On dérive :
\[
g'(x)=2x+2.
\]
Le point critique de $g$ est donné par
\[
2x+2=0
\qquad\Longleftrightarrow\qquad
x=-1.
\]
Comme $-\sqrt{2}<-1<\sqrt{2}$, ce point appartient bien à l’intervalle.

Le signe de $g'(x)$ est :
\begin{itemize}
\item négatif si $x<-1$ ;
\item nul si $x=-1$ ;
\item positif si $x>-1$.
\end{itemize}
Donc $g$ est décroissante sur $[-\sqrt{2},-1]$, puis croissante sur $[-1,\sqrt{2}]$.

Valeurs remarquables :
\[
g(-\sqrt{2})=2-2\sqrt{2},
\qquad
g(-1)=1-2=-1,
\qquad
g(\sqrt{2})=2+2\sqrt{2}.
\]

Ainsi, sur $B_1$ :
\begin{itemize}
\item le minimum est $-1$, atteint en $(-1,-1)$ ;
\item le maximum est $2+2\sqrt{2}$, atteint en $(\sqrt{2},\sqrt{2})$.
\end{itemize}
}

\item   \question{Étudier la restriction de $f$ à $B_2$. On posera
$
h(\theta)=f(2\cos\theta,2\sin\theta).
$
Déterminer les points critiques de $h$ sur $[\tfrac{\pi}{4},\tfrac{5\pi}{4}]$, puis ses extremums sur cet intervalle.}
\indication{}
\reponse{
Sur $B_2$, on pose
\[
x=2\cos\theta,\qquad y=2\sin\theta,
\qquad \theta\in\left[\frac{\pi}{4},\frac{5\pi}{4}\right].
\]
Alors
\[
h(\theta)=f(2\cos\theta,2\sin\theta)
=(2\cos\theta)^2+2(2\sin\theta)
=4\cos^2\theta+4\sin\theta.
\]

On dérive :
\[
h'(\theta)=-8\sin\theta\cos\theta+4\cos\theta
=4\cos\theta(1-2\sin\theta).
\]
Les points critiques vérifient
\[
4\cos\theta(1-2\sin\theta)=0.
\]
Donc :
\[
\cos\theta=0
\qquad\text{ou}\qquad
\sin\theta=\frac12.
\]

Sur l’intervalle $\left[\frac{\pi}{4},\frac{5\pi}{4}\right]$, cela donne :
\[
\theta=\frac{\pi}{2}
\qquad\text{et}\qquad
\theta=\frac{5\pi}{6}.
\]

Étudions ensuite les valeurs de $h$ aux points critiques et aux bornes :
\[
h\left(\frac{\pi}{4}\right)=4\cdot \frac12+4\cdot \frac{\sqrt{2}}{2}=2+2\sqrt{2},
\]
\[
h\left(\frac{\pi}{2}\right)=4\cdot 0+4\cdot 1=4,
\]
\[
h\left(\frac{5\pi}{6}\right)=4\cdot \frac34+4\cdot \frac12=3+2=5,
\]
\[
h\left(\frac{5\pi}{4}\right)=4\cdot \frac12+4\cdot \left(-\frac{\sqrt{2}}{2}\right)=2-2\sqrt{2}.
\]

On peut aussi étudier le signe de $h'$ :
\begin{itemize}
\item sur $\left[\frac{\pi}{4},\frac{\pi}{2}\right]$, $h$ décroît ;
\item sur $\left[\frac{\pi}{2},\frac{5\pi}{6}\right]$, $h$ croît ;
\item sur $\left[\frac{5\pi}{6},\frac{5\pi}{4}\right]$, $h$ décroît.
\end{itemize}

Ainsi, sur $B_2$ :
\begin{itemize}
\item le minimum est
\[
2-2\sqrt{2},
\]
atteint pour
\[
\theta=\frac{5\pi}{4},
\]
c’est-à-dire au point
\[
(-\sqrt{2},-\sqrt{2}) ;
\]
\item le maximum est
\[
5,
\]
atteint pour
\[
\theta=\frac{5\pi}{6},
\]
c’est-à-dire au point
\[
\left(2\cos\frac{5\pi}{6},2\sin\frac{5\pi}{6}\right)=(-\sqrt{3},1).
\]
\end{itemize}
}

\item   \question{En déduire le minimum global et le maximum global de $f$ sur $K$, ainsi que les points où ils sont atteints.}
\indication{}
\reponse{
Il n’y a pas de point critique dans l’intérieur de $K$, donc les extrema globaux sont atteints sur le bord, c’est-à-dire sur $B_1\cup B_2$.

D’après les questions précédentes :
\begin{itemize}
\item sur $B_1$, le minimum vaut $-1$ en $(-1,-1)$, et le maximum vaut $2+2\sqrt{2}$ en $(\sqrt{2},\sqrt{2})$ ;
\item sur $B_2$, le minimum vaut $2-2\sqrt{2}$ en $(-\sqrt{2},-\sqrt{2})$, et le maximum vaut $5$ en $(-\sqrt{3},1)$.
\end{itemize}

On compare ces valeurs :
\[
2-2\sqrt{2}\approx -0{,}828,
\qquad
-1<2-2\sqrt{2},
\]
donc le minimum global est
\[
-1,
\]
atteint au point
\[
(-1,-1).
\]

De même,
\[
5>2+2\sqrt{2}\approx 4{,}828,
\]
donc le maximum global est
\[
5,
\]
atteint au point
\[
(-\sqrt{3},1).
\]

Conclusion :
\[
\min_K f=-1 \quad \text{atteint en } (-1,-1),
\]
\[
\max_K f=5 \quad \text{atteint en } (-\sqrt{3},1).
\]
}

\end{enumerate}
\vspace{0.5cm}
}